\documentclass[10pt,a4paper]{amsart}
\usepackage[margin=19mm]{geometry}
\usepackage{amsmath,amssymb,mathtools}
\usepackage[scr=rsfs,cal=euler]{mathalfa}
\usepackage{xfrac}
\usepackage[unicode,hidelinks,bookmarks=false]{hyperref}
\newcommand{\ie}{i.e.\ }
\newcommand{\cf}{cf.\ }
\newcommand{\resp}{respectively\ }
\newcommand{\Subsection}{\subsection}
\providecommand{\nobreakdash}{\nobreak\hbox{-}}
\setlength{\emergencystretch}{2em}
\multlinegap0pt
\usepackage{amssymb}
\usepackage{algorithm}\usepackage{algpseudocode}
\usepackage{enumerate}
\newtheorem{lemma}{Lemma}
\title{Cyclic Relaxed Douglas-Rachford Splitting for Inconsistent Nonconvex Feasibility}
\author{Thi Lan Dinh, G. S. Matthijs Jansen, D. Russell Luke}
\date{Source: arXiv:2502.12285v1 (CC BY 4.0)}
\begin{document}
\maketitle

\noindent\emph{Source and licence.} Cyclic Relaxed Douglas-Rachford Splitting for Inconsistent Nonconvex Feasibility. Version arXiv:2502.12285v1 (CC BY 4.0), distributed by arXiv under the Creative Commons Attribution 4.0 International licence, http://creativecommons.org/licenses/by/4.0/, as recorded in the arXiv licence metadata for this exact version and shown on the abstract page https://arxiv.org/abs/2502.12285v1. Full licence notice: \texttt{licenses/arxiv-2502.12285-CC-BY-4.0.txt}.\par\smallskip

\noindent\textit{Editorial scope.} Counted unit: the article's lemma lem:composition.affine.2 (source0.tex lines 925--937). The article's complete original proof of it is delivered with it (source0.tex lines 938--957, one \texttt{proof} environment of 20 lines). No mathematics is written, repaired or re-derived: the statement and the proof are the article's own lines, in the article's own notation, and no retained line is substituted or re-flowed.  No further source-local context is needed: every cross-reference of every retained line resolves inside the two delivered spans.  Nothing was omitted from any retained span, so the package carries no omission record.  No retained line contains a citation, so the wrapper carries no bibliography and no bibliography entry.  No retained line contains a figure environment, an image-inclusion command or drawing code, so no image is delivered and the package declares no asset dependency.  Editorial formatting only, in addition: an AMS \texttt{amsart} preamble that restates the article's own declarations and macros used by the retained lines, namely the environments \texttt{lemma} on separate counters, together with the standard packages those lines need.  The retained environments print in this document as Lemma 1 in order of appearance, whereas the article prints the same items with the numbering of its own full document.  The complete native source of the article as distributed in the arXiv e-print for this exact version is bundled unchanged as \texttt{sources/173441/source0.tex}.
\begin{lemma}\label{lem:composition.affine.2}
Let $T_i:\mathbb E\rightrightarrows \mathbb E$ be an operator that is approximated by an affine operator $\Phi_i:\mathbb E\to \mathbb E$ on a set $U_i\subset \mathbb E$ with error $\varepsilon_i>0$,  $i=1,2$.
Assume that $T_2U_2\subseteq U_1$.
Then $T_1\circ T_2$ is approximated by the affine operator $\Phi_{1,2}$ on $U_2$ with
\begin{equation}
\Phi_{1,2}=(1-\varepsilon_2)\Phi_1(\frac{1}{1-\varepsilon_2}\Phi_2)
\end{equation}
and error
\begin{equation}
\left\|T_1\circ T_2x- \Phi_{1,2}x \right\|\leq
\varepsilon_{1,2}:=\varepsilon_1+\varepsilon_2 \sup_{w\in\mathbb B}{\|\Phi_1w\|}\quad \forall x\in U_2.
\end{equation}
\end{lemma}

\begin{proof}
Let $x\in U_2$.
Let $y\in T_2x$.
By assumption, $y=\Phi_2 x +\varepsilon_2 v_2$ for some $v_2\in\mathbb B$.
Since $T_2U_2\subseteq U_1$, we get $y\in U_1$.
Let $z\in T_1y$.
By assumption, there exists $v_1\in \mathbb B$ such that
\begin{equation}
\begin{array}{rl}
z=&\Phi_1 y+\varepsilon_1 v_1\\
=&\Phi_1 (\Phi_2 x +\varepsilon_2 v_2)+\varepsilon_1 v_1\\
=&\Phi_1 ((1-\varepsilon_2)\frac{1}{1-\varepsilon_2}\Phi_2 x +\varepsilon_2 v_2)+\varepsilon_1 v_1\\
=&(1-\varepsilon_2)\Phi_1 (\frac{1}{1-\varepsilon_2}\Phi_2 x) +\varepsilon_2 \Phi_1v_2+\varepsilon_1 v_1\\
=&\Phi_{1,2} x +\varepsilon_{1,2} v\,,
\end{array}
\end{equation}
where $v=\frac{1}{\varepsilon_{1,2}}(\varepsilon_2 \Phi_1v_2+\varepsilon_1 v_1)$.
Note that $\|v\|\le 1$.
Hence the result follows.
\end{proof}

\end{document}
